\input{../tex-inputs/latex-leseansicht-vorspann}

\section[Arthur und Olga Schnitzler an Gustav Schwarzkopf, 14. 1. 1903]{L04525 Arthur und Olga Schnitzler an Gustav Schwarzkopf, 14. 1. 1903}
\nopagebreak\mylabel{L04525v}
\rehead{ }\normalsize\beginnumbering\briefempfaengerindex{Schwarzkopf, Gustav@\textsc{Schwarzkopf, Gustav}!zzzSchnitzler, Olga@\emph{von Olga Schnitzler}!1903-01-142@{14. 1. 1903}|(be}\briefempfaengerindex{Schwarzkopf, Gustav@\textsc{Schwarzkopf, Gustav}!zzzSchnitzler, Arthur@\emph{von Arthur Schnitzler}!1903-01-142@{14. 1. 1903}|(be}
\toendnotes[C]{\smallbreak\pagebreak[2]}
\correspDesc{Versand  durch Arthur Schnitzler, Olga Schnitzler am 14. 1. 1903 in Salzburg
\newline{}Erhalt  durch Gustav Schwarzkopf im Zeitraum [15. 1. 1903 – 19. 1. 1903?] in Wien}\toendnotes[C]{\smallbreak}
\Standort{DLA, A:Schnitzler, HS.1985.1.1897.}
\physDesc{Bildpostkarte, 150 Zeichen
\newline{}Handschrift Arthur Schnitzler: Bleistift, deutsche Kurrent
\newline{}Handschrift Olga Schnitzler: Bleistift, lateinische Kurrent
\newline{}Versand: 1) Stempel: »\nobreak{}\oindex{Salzburg@\textbf{Salzburg}|pwk}Sa\textcolor{gray}{lzburg}, 6 – \textcolor{gray}{7 N}\nobreak{}«.   2) Stempel: »\nobreak{}\oindex{I., Innere Stadt@\textbf{I., Innere Stadt}|pwk}Wien 1\textcolor{gray}{/1}, 15. 1. 03, 8 – \textcolor{gray}{9 V.}, 1, Bestellt\nobreak{}«. }
\pstart
           \noindent{}\centering{}{\pb}\textcolor{gray}{\textbf{Neuthor\oindex{Sigmundstor@\textbf{Sigmundstor}|pw}}}\hspace*{2.5em}\textcolor{gray}{\textbf{Salzburg\oindex{Salzburg@\textbf{Salzburg}|pw}}}\pend
           \pstart{}{\pb}Herrn Guſtav Schwarzkopf\pend{}\pstart{}Wien I\oindex{I., Innere Stadt@\textbf{I., Innere Stadt}|pw}\pend{}\pstart{}Tiefer Graben 23\oindex{Wien@\textbf{Wien}!I., Innere Stadt@\textbf{I., Innere Stadt}!Tiefer Graben 23@\textbf{Tiefer Graben 23}|pw}.\pend{}{\bigskip}\vspace{1em}
\pstart
           \centering{}14.\,1.\,903\pend
           \vspace{0.5em}
\pstart
           Wir grüßen Sie herzlich, Dr. Max\pwindex{Schwarzkopf, Max 12.\,6.\,1857 Wien – 14.\,4.\,1928 ebd.@\textsc{Schwarzkopf, Max} (12.\,6.\,1857 Wien – 14.\,4.\,1928 ebd.)|pw} desgleichen,\pend
           \pstart \spacefill\mbox{A.}\pend{}\selectlanguage{ngerman}\vspace{1em}
\pstart
           \noindent{}{[}hs. O. Schnitzler:{]} Grüß Sie Gott!\pend
           \pstart \spacefill\mbox{Olga}\pend{}\selectlanguage{ngerman}\vspace{1em}
\pstart
           \noindent{}{[}hs. A. Schnitzler:{]} So ſchön und weiſs und ſtill!\pend
           \selectlanguage{ngerman}\endnumbering\briefempfaengerindex{Schwarzkopf, Gustav@\textsc{Schwarzkopf, Gustav}!zzzSchnitzler, Olga@\emph{von Olga Schnitzler}!1903-01-142@{14. 1. 1903}|)be}\briefempfaengerindex{Schwarzkopf, Gustav@\textsc{Schwarzkopf, Gustav}!zzzSchnitzler, Arthur@\emph{von Arthur Schnitzler}!1903-01-142@{14. 1. 1903}|)be}\mylabel{L04525h}
\begin{anhang}
\end{anhang}\newcommand{\dateiname}{L04525}\newcommand{\titel}{Arthur und Olga Schnitzler an Gustav Schwarzkopf, 14. 1. 1903}\newcommand{\editorInnen}{Herausgegeben von Jahnke, SelmaMüller, Martin Anton}\input{../tex-inputs/latex-leseansicht-abspann}
